\documentclass[11pt]{article}

\usepackage[left=1.5cm, right=1.5cm, top=2cm, bottom=2cm]{geometry}
\usepackage{amsmath}
\usepackage{amssymb}
\usepackage{amsfonts}
\usepackage{physics}
\usepackage{booktabs}
\usepackage{graphicx}
\usepackage{bm}
\usepackage{algorithm}
\usepackage{algpseudocode}

\setcounter{MaxMatrixCols}{20}

\title{Kanamori Note}
\author{Anderson Steckler}
\date{}
\begin{document}
\maketitle
\section{Hamiltonian}
Our single layer Hamiltonian is given by the following:
\begin{gather}
    H = H_0 + H_{SOC} + H_{Kanamori}
\end{gather}
the tight binding part of the Hamiltonian is given by (yz, xz, xy) orbital basis.
\begin{gather}
    e_{xy} = -2 t_1(\cos k_x + \cos k_y)-4 t_2 \cos k_x \cos k_y \\
    e_{yz} = -2 t_1 \cos k_x - 2 t_{\delta} \cos k_y \\
    e_{zx} = -2 t_\delta \cos k_x - 2 t_1 \cos k_y\\
    \\
    H_0(k) = \begin{pmatrix}
        e_{yz} & 0 & 0 \\
        0 & e_{xz} & 0 \\
        0 & 0 & e_{xy} + \delta_{CF}
    \end{pmatrix}
\end{gather}
where $\Delta_{CF}$ is the crystal field splitting. For the hoppings $t_1$ represents 
the in plane $V_{dd\pi}$ overlap. $t_2$ represents the diagonal nearest neighbor hopping.
$t_\delta$ represents a weak hopping which comes from the $V_{dd\delta}$ overlap.
The spin orbit coupling term is given by
\begin{gather}
    H_{SOC} = \lambda \bm{L} \cdot \bm{S}
\end{gather}
where $\lambda$ is the strength of the spin orbit coupling. 
The Kanamori interaction term is done in mean-field and derivation is given at the end. The result is this:
\begin{align}
    H^k_{MF} &= U \sum_m \langle n_{m\uparrow}\rangle \hat n_{m\downarrow} + \langle n_{m\downarrow}\rangle \hat n_{m\uparrow} \\\nonumber
    &+ U' \sum_{m \neq m'} \langle n_{m\uparrow}\rangle \hat n_{m'\downarrow}
    + \langle n_{m' \downarrow}\rangle \hat n_{m\uparrow}  \\\nonumber
    &+ (U'-J) \sum_{m < m', \sigma} \langle n_{m\sigma} \rangle \hat n_{m'\sigma}
    + \langle n_{m' \sigma} \rangle \hat n_{m\sigma} \\\nonumber
    &+ J\sum_{m\neq m'}\left[\langle d^\dagger_{m\uparrow} d_{m'\uparrow} \rangle d^\dagger_{m' \downarrow} d_{m\downarrow}
    + \langle d^\dagger_{m' \downarrow} d_{m\downarrow}\rangle d^\dagger_{m\uparrow}d_{m' \uparrow}
    - \langle d^\dagger_{m \uparrow} d_{m\downarrow}\rangle d^\dagger_{m'\downarrow}d_{m'\uparrow}
    - \langle d^\dagger_{m'\downarrow} d_{m'\uparrow} \rangle d^\dagger_{m\uparrow} d_{m \downarrow}\right] \\\nonumber
    &+ J\sum_{m\neq m'}\left[
        \langle d^\dagger_{m\uparrow} d_{m'\uparrow}\rangle d^\dagger_{m\downarrow} d_{m'\downarrow}
        +\langle d^\dagger_{m\downarrow} d_{m'\downarrow}\rangle d^\dagger_{m\uparrow}d_{m'\uparrow}
        - \langle d^\dagger_{m\uparrow} d_{m'\downarrow}\rangle d^\dagger_{m\downarrow}d_{m'\uparrow}
        - \langle d^\dagger_{m\downarrow} d_{m'\uparrow}\rangle d^\dagger_{m\uparrow} d_{m'\downarrow}
    \right]
\end{align}

\section{Bilayer Hamiltonian}
The bilayer Hamiltonian is given by
\begin{gather}
    H = \begin{pmatrix}
        H_0 & H_\perp \\
        H_\perp & H_0
    \end{pmatrix}
\end{gather}
With the interlayer hoppings given by
\begin{gather}
    H_\perp = \begin{pmatrix}
        t_{\perp} & 0 & 0 \\
        0 & t_{\perp} & 0 \\
        0 & 0 & t_{xy}
    \end{pmatrix}
\end{gather}

\section{Basis Ordering}
Our order is layer-\> spin -\> orbital. 
\begin{table}[h]
\centering
\begin{tabular}{cccc}
\hline
\textbf{Index} & \textbf{Layer} & \textbf{Spin} & \textbf{Orbital} \\
\hline
0  & $L_1$ & $\uparrow$ & $d_{yz}$ \\
1  & $L_1$ & $\uparrow$ & $d_{xz}$ \\
2  & $L_1$ & $\uparrow$ & $d_{xy}$ \\
3  & $L_1$ & $\downarrow$ & $d_{yz}$ \\
4  & $L_1$ & $\downarrow$ & $d_{xz}$ \\
5  & $L_1$ & $\downarrow$ & $d_{xy}$ \\
6  & $L_2$ & $\uparrow$ & $d_{yz}$ \\
7  & $L_2$ & $\uparrow$ & $d_{xz}$ \\
8  & $L_2$ & $\uparrow$ & $d_{xy}$ \\
9  & $L_2$ & $\downarrow$ & $d_{yz}$ \\
10 & $L_2$ & $\downarrow$ & $d_{xz}$ \\
11 & $L_2$ & $\downarrow$ & $d_{xy}$ \\
\hline
\end{tabular}
\caption{Basis state ordering: layer-major $\to$ spin-major $\to$ orbital-minor.}
\label{tab:basis_ordering}
\end{table}

\section{Self Consistent Loop}
Our order parameter is the density matrix
\begin{gather}
    \rho_{ab} = \langle c^\dagger_a c_b \rangle = \frac{1}{N_k} 
    \sum_{\mathbf{k},n} f(\varepsilon_{n\mathbf{k}} - \mu)\, v_{n\mathbf{k},a}^* \, v_{n\mathbf{k},b}
\end{gather}
where a,b run over the 12 states. Wit our parameters we first run the SCF loop with 
just the hubbard hamiltonian to get a diagonal initial guess for the density matrix.
Then we break the symmetry. Doing this in different ways can converge on different
results so you would have to take the one of lowest energy. For example I get interesting
results comparable to the DFT with this.
\begin{equation*}
\delta\rho = \delta
\begin{pmatrix}
    +1 & +1 &  0 &  0 &  0 &  0 &  0 &  0 &  0 &  0 &  0 &  0 \\
    +1 &  0 &  0 &  0 &  0 &  0 &  0 &  0 &  0 &  0 &  0 &  0 \\
    0 &  0 & -1 &  0 &  0 &  0 &  0 &  0 &  0 &  0 &  0 &  0 \\
    0 &  0 &  0 & +1 &  0 &  0 &  0 &  0 &  0 &  0 &  0 &  0 \\
    0 &  0 &  0 &  0 & +1 &  0 &  0 &  0 &  0 &  0 &  0 &  0 \\
    0 &  0 &  0 &  0 &  0 & +1 &  0 &  0 &  0 &  0 &  0 &  0 \\
    0 &  0 &  0 &  0 &  0 &  0 & +1 & +1 &  0 &  0 &  0 &  0 \\
    0 &  0 &  0 &  0 &  0 &  0 & +1 &  0 &  0 &  0 &  0 &  0 \\
    0 &  0 &  0 &  0 &  0 &  0 &  0 &  0 & -1 &  0 &  0 &  0 \\
    0 &  0 &  0 &  0 &  0 &  0 &  0 &  0 &  0 & -1 &  0 &  0 \\
    0 &  0 &  0 &  0 &  0 &  0 &  0 &  0 &  0 &  0 & -1 &  0 \\
    0 &  0 &  0 &  0 &  0 &  0 &  0 &  0 &  0 &  0 &  0 & -1 \\
\end{pmatrix}
\end{equation*}
Which AI reccomended I need to think about the meaning of each of these and why. We then
do the diagonalization and calculate the new density matrix using the convergance criteria as
\begin{gather}
    || \rho_{calc} - \rho_{old} || < \epsilon = 10^{-6}
\end{gather}
where this is the Frobenius norm. Flatten the matrix and take the norm. The whole process can be summarized
\begin{algorithm}[h]
\caption{Kanamori Self-Consistent Field}
\label{alg:kanamori_scf}
\begin{algorithmic}[1]
\State $\rho \leftarrow \rho_0$ \Comment{Initial guess from converged Stoner SCF}
\For{each iteration}
    \State Build $H(\mathbf{k})$ from $\rho$ via \texttt{compute\_eigensystem\_kanamori}
    \State Find $\mu$ via Brent's method on $N(\mu) = N_{\text{target}}$
    \State Compute $\rho_{\text{new}}$ from eigensystem and Fermi-Dirac weights
    \If{$\|\rho_{\text{new}} - \rho\|_F < 10^{-6}$}
        \State Compute total energy and \Return
    \EndIf
    \State Mix: $\rho \leftarrow \alpha\, \rho_{\text{new}} + (1-\alpha)\, \rho$
\EndFor
\end{algorithmic}
\end{algorithm}

the alpha I have been using is $\alpha = 0.2$.

\section{Results}
\subsection{Hunds Rule}
We turn off hopping and SOC. With $U=4~eV$ and $J=0.8~eV$. We also set $t_\perp$ to zero
so the layers are decoupled.
\begin{table}[h]
\centering
\begin{tabular}{lcc}
\hline
Orbital & Spin Up & Spin Down \\
\hline
$d_{xy}$ & 1.0000 & 0.6667 \\
$d_{xz}$ & 1.0000 & 0.6667 \\
$d_{yz}$ & 1.0000 & 0.6667 \\
\hline
Total & 5.0000 & \\
\hline
\end{tabular}
\caption{Orbital occupancies.}
\label{tab:occupancy}
\end{table}
this follows the first Hund's rule of maxamizing S. But does not maximize L. This is because
the guess for $\rho_0$ was diagonal. Adding some shift (I didn't write down exactly what it was)
we get
\begin{table}[h]
\centering
\begin{tabular}{lcc}
\hline
Orbital & Spin Up & Spin Down \\
\hline
dxy & 1.0000 & 1.0000 \\
dxz & 1.0000 & 1.0000 \\
dyz & 1.0000 & 0.0000 \\
\hline
Total & 5.0000 & \\
\hline
\end{tabular}
\caption{Orbital occupancies.}
\label{tab:occupancy}
\end{table}
Where now the yz is unoccupied which corresponds to angular of $m=-1$ in the TP equivelence.
So now the second rule is satisfied of maximizing L.
\\\\
One more comment is that I tried some different shift and got that my code did not converge.
Again I forgot to write it down which is too bad, but I can try to reproduce it later.
I think in general the shift to the diagonal $\rho_0$ needs to align with the expected result.

\subsection{Charge Disportion}
I added a different splitting of the $xy$ orbital for each layer, one positive and one negative
with slightly less absoulte value. Using the shift to the initial density matrix above (the big 12x12) one
I get.
\begin{table}[h]
\centering
\begin{tabular}{lccc}
\hline
 & $n_{xz/yz\downarrow}$ & $n_{xy\downarrow}$ & $n_{\downarrow}$ (total t2g) \\
\hline
Co1 DFT   & 0.98 & 0.32 & 2.28 \\
Co1 Model & 0.78 & 0.37 & 1.94 \\
Co2 DFT   & 0.38 & 1.00 & 1.76 \\
Co2 Model & 0.62 & 0.82 & 2.06 \\
\hline
\end{tabular}
\caption{Comparison of Sr3Co2O7 DFT (table 2) and model orbital occupancies.}
\label{tab:comparison}
\end{table}
If we sum the occupations and subtract each layer we get a similar value to the dft $0.13$ vs $0.14$,
so these crystal field values capture this result. Likewise we see in one layer the $xy$
orbital has the higher occupation which is reverse to the other layer. But the ratios are not exactly the same.
This could be due to a mean field artifact, due to it following Hund's first rule it wants the occupations
of the spin up channel to always be one. 
\\\\
We used the following parameters.
\begin{table}[h]
\centering
\begin{tabular}{lc}
\hline
Parameter & Value \\
\hline
$t_1$              & 0.3 eV   \\
$t_\delta$         & 0.05 eV  \\
$t_2$              & 0.03 eV  \\
$t_\perp$          & 0.2 eV   \\
$t_\perp^{xy}$     & 0.05 eV  \\
$\lambda$          & 0.1 eV   \\
$U$                & 4.0 eV   \\
$J$                & 0.8 eV   \\
$U'$               & 2.4 eV   \\
$\delta_{cf}^{L1}$ & +0.23 eV \\
$\delta_{cf}^{L2}$ & -0.16 eV \\
\hline
\end{tabular}
\caption{Model parameters.}
\label{tab:parameters}
\end{table}

These were just suggested from AI I need to go through on my own and think about
if these are reasonable.
\\\\
These are the shared parameters between the Initial Hubabrd Loop and the Kanamori loop.
\begin{table}[h]
\centering
\begin{tabular}{lll}
\hline
Parameter & Value & Description \\
\hline
$S_0$       & 0.3  & Initial magnetization \\
$\alpha$    & 0.2  & Mixing parameter \\
$T$         & 0.05 eV & Temperature \\
$N_{\text{target}}$ & 10.0 & Target electron number (5 per layer $\times$ 2 layers) \\
grid        & 50   & $k$-point grid size \\
\hline
\end{tabular}
\caption{SCF numerical parameters.}
\label{tab:scf_params}
\end{table}

$S0 = m/2$ and is the initial magnetization for hubbard loop.

\section{Kanamori Hamiltonian}
To capture the complete physics we extend out model to the Kanamori Hamiltonian
\begin{align}
    H_K &= U\sum_m \hat n_{m \uparrow} \hat n_{m \downarrow} 
    + U' \sum_{m \neq m'} \hat n_{m \uparrow}  \hat n_{m' \downarrow} 
    + (U'-J) \sum_{m < m', \sigma} \hat n_{m\sigma} \hat n_{m' \sigma} \\\nonumber
    &- J \sum_{m \neq m'} d^\dagger_{m\uparrow} d_{m \downarrow}d^\dagger_{m' \downarrow}d_{m'\uparrow}
    + J \sum_{m \neq m'} d^\dagger_{m\uparrow} d^\dagger_{m \downarrow} d_{m' \downarrow} d_{m' \uparrow}
\end{align}
We have already looked at the MF treatment of the first term above.

\subsection{Onsite Repulsion}
\begin{gather}
        H = U \sum_m \langle n_{m\uparrow}\rangle \hat n_{m\downarrow} + \langle n_{m\downarrow}\rangle \hat n_{m\uparrow} - \langle n_{m\uparrow}\rangle \langle n_{m\downarrow}\rangle
\end{gather}
\subsection{Interorbital Repulsion}
\subsubsection{Opposite Spin}
\begin{gather}
    H = U' \sum_{m \neq m'} \hat n_{m \uparrow}  \hat n_{m' \downarrow} 
\end{gather}
As before each operator is written as
\begin{gather}
    \hat n_{m\sigma} = \langle n_{m\sigma} \rangle + \delta n_{m \sigma}
\end{gather}
which allows us to write
\begin{align}
    H &= U' \sum_{m\neq m'} [\langle n_{m\uparrow} \rangle + \delta n_{m \uparrow}]
    [\langle n_{m'\downarrow} \rangle + \delta n_{m' \downarrow}] \\\nonumber
    &= U' \sum_{m \neq m'} \langle n_{m \uparrow}\rangle \langle n_{m' \downarrow}\rangle 
    + \langle n_{m\uparrow}\rangle \delta n_{m' \downarrow} + \langle n_{m' \downarrow}\rangle \delta n_{m \uparrow}
\end{align}
with $\delta n_{m\sigma} = \hat n_{m\sigma} - \langle n_{m\sigma}\rangle$
\begin{align}
    H &= U' \sum_{m\neq m'} \langle n_{m \uparrow}\rangle \langle n_{m' \downarrow}\rangle
    + \langle n_{m \uparrow}\rangle [\hat n_{m' \downarrow} - \langle n_{m'\downarrow}\rangle]
    + \langle n_{m'\downarrow}\rangle [\hat n_{m\uparrow} - \langle n_{m\uparrow}\rangle] \\\nonumber
    &= U' \sum_{m \neq m'} \langle n_{m\uparrow}\rangle \hat n_{m'\downarrow}
    + \langle n_{m' \downarrow}\rangle \hat n_{m\uparrow} 
    -\langle n_{m \uparrow}\rangle \langle n_{m' \downarrow}\rangle
\end{align}
\subsubsection{Same Spin}
The next term is
\begin{align}
    H &= (U'-J) \sum_{m < m', \sigma} \hat n_{m\sigma} \hat n_{m' \sigma} \\\nonumber
      &= (U'-J) \sum_{m < m', \sigma} [\langle n_{m \sigma}\rangle + \delta n_{m\sigma}][\langle n_{m'\sigma}\rangle + \delta n_{m'\sigma}]\\\nonumber
      &= (U'-J) \sum_{m < m', \sigma} \langle n_{m\sigma}\rangle\langle n_{m'\sigma}\rangle
      + \langle n_{m\sigma}\rangle \delta n_{m' \sigma} + \langle n_{m'\sigma}\rangle\delta n_{m\sigma}
\end{align}
with $\delta n_{m\sigma} = \hat n_{m\sigma} - \langle n_{m\sigma}\rangle$
\begin{align}
    H &= (U'-J) \sum_{m < m', \sigma} \langle n_{m\sigma}\rangle\langle n_{m'\sigma}\rangle
      + \langle n_{m\sigma}\rangle [\hat n_{m'\sigma} - \langle n_{m'\sigma}\rangle] 
      + \langle n_{m'\sigma}\rangle[\hat n_{m\sigma} - \langle n_{m\sigma}\rangle] \\\nonumber
      &= (U'-J) \sum_{m < m', \sigma} \langle n_{m\sigma} \rangle \hat n_{m'\sigma}
      + \langle n_{m' \sigma} \rangle \hat n_{m \sigma}
      - \langle n_{m\sigma} \rangle \langle n_{m' \sigma} \rangle
\end{align}
\subsection{Exchange/Pair-Hopping Term}
The terms are
\begin{align}
    H &= -J \sum_{m \neq m'} d^\dagger_{m\uparrow} d_{m \downarrow}d^\dagger_{m' \downarrow}d_{m'\uparrow}
    + J \sum_{m \neq m'} d^\dagger_{m\uparrow} d^\dagger_{m \downarrow} d_{m' \downarrow} d_{m' \uparrow}
\end{align}
We can normal order and get
\begin{align}
    H &= J \sum_{m \neq m'} \left[d^\dagger_{m\uparrow} d^\dagger_{m' \downarrow}d_{m \downarrow}d_{m'\uparrow}
    + d^\dagger_{m\uparrow} d^\dagger_{m \downarrow} d_{m' \downarrow} d_{m' \uparrow}\right]
\end{align}
Here we expand with mean field
\begin{align}
    d^\dagger_1 d^\dagger_2 d_3 d_4 = 
    -\langle d^\dagger_1 d_3 \rangle d^\dagger_2 d_4
    +\langle d^\dagger_1 d_4 \rangle d_2^\dagger d_3
    -\langle d^\dagger_2 d_4 \rangle d^\dagger_1 d_3
    + \langle d^\dagger_2 d_3 \rangle d^\dagger_1 d_4 \\\nonumber
    + DC
\end{align}
Where DC corrects for double counting and is
\begin{gather}
    DC = \langle d^\dagger_1 d_4\rangle\langle d^\dagger_2 d_3 \rangle 
    - \langle d^\dagger_1 d_3 \rangle \langle d^\dagger_2 d_4 \rangle
\end{gather}
Where we drop the $\langle d d \rangle$ and the $\langle d^\dagger d^\dagger \rangle$ channels.
Hence we can write
\begin{gather}
    H_{exc} = J\sum_{m\neq m'}\left[\langle d^\dagger_{m\uparrow} d_{m'\uparrow} \rangle d^\dagger_{m' \downarrow} d_{m\downarrow}
    + \langle d^\dagger_{m' \downarrow} d_{m\downarrow}\rangle d^\dagger_{m\uparrow}d_{m' \uparrow}
    - \langle d^\dagger_{m \uparrow} d_{m\downarrow}\rangle d^\dagger_{m'\downarrow}d_{m'\uparrow}
    - \langle d^\dagger_{m'\downarrow} d_{m'\uparrow} \rangle d^\dagger_{m\uparrow} d_{m \downarrow}\right]
\end{gather}
with a correction to the energy shift given by the following
\begin{gather}
    \Delta E_{exc} = J\sum_{m\neq m'} \left[
        \langle d^\dagger_{m\uparrow} d_{m'\uparrow}\rangle \langle d^\dagger_{m'\downarrow} d_{m\downarrow}\rangle
        -\langle d^\dagger_{m\uparrow} d_{m\downarrow}\rangle \langle d^\dagger_{m'\downarrow} d_{m'\uparrow}\rangle
    \right]
\end{gather}
Likewise we can write
\begin{gather}
    H_{ph} = J\sum_{m\neq m'}\left[
        \langle d^\dagger_{m\uparrow} d_{m'\uparrow}\rangle d^\dagger_{m\downarrow} d_{m'\downarrow}
        +\langle d^\dagger_{m\downarrow} d_{m'\downarrow}\rangle d^\dagger_{m\uparrow}d_{m'\uparrow}
        - \langle d^\dagger_{m\uparrow} d_{m'\downarrow}\rangle d^\dagger_{m\downarrow}d_{m'\uparrow}
        - \langle d^\dagger_{m\downarrow} d_{m'\uparrow}\rangle d^\dagger_{m\uparrow} d_{m'\downarrow}
    \right]
\end{gather}
\begin{gather}
    \Delta E_{ph} = J \sum_{m\neq m'}\left[ \langle d^\dagger_{m\uparrow} d_{m'\uparrow} \rangle 
    \langle d^\dagger_{m\downarrow} d_{m'\downarrow} \rangle
    - \langle d^\dagger_{m\uparrow} d_{m' \downarrow} \rangle
    \langle d^\dagger_{m\downarrow}d_{m'\uparrow}\rangle \right]
\end{gather}
\section{Complete Kanamori Hamiltonian}
\begin{align}
    H^k_{MF} &= U \sum_m \langle n_{m\uparrow}\rangle \hat n_{m\downarrow} + \langle n_{m\downarrow}\rangle \hat n_{m\uparrow} \\\nonumber
    &+ U' \sum_{m \neq m'} \langle n_{m\uparrow}\rangle \hat n_{m'\downarrow}
    + \langle n_{m' \downarrow}\rangle \hat n_{m\uparrow}  \\\nonumber
    &+ (U'-J) \sum_{m < m', \sigma} \langle n_{m\sigma} \rangle \hat n_{m'\sigma}
    + \langle n_{m' \sigma} \rangle \hat n_{m\sigma} \\\nonumber
    &+ J\sum_{m\neq m'}\left[\langle d^\dagger_{m\uparrow} d_{m'\uparrow} \rangle d^\dagger_{m' \downarrow} d_{m\downarrow}
    + \langle d^\dagger_{m' \downarrow} d_{m\downarrow}\rangle d^\dagger_{m\uparrow}d_{m' \uparrow}
    - \langle d^\dagger_{m \uparrow} d_{m\downarrow}\rangle d^\dagger_{m'\downarrow}d_{m'\uparrow}
    - \langle d^\dagger_{m'\downarrow} d_{m'\uparrow} \rangle d^\dagger_{m\uparrow} d_{m \downarrow}\right] \\\nonumber
    &+ J\sum_{m\neq m'}\left[
        \langle d^\dagger_{m\uparrow} d_{m'\uparrow}\rangle d^\dagger_{m\downarrow} d_{m'\downarrow}
        +\langle d^\dagger_{m\downarrow} d_{m'\downarrow}\rangle d^\dagger_{m\uparrow}d_{m'\uparrow}
        - \langle d^\dagger_{m\uparrow} d_{m'\downarrow}\rangle d^\dagger_{m\downarrow}d_{m'\uparrow}
        - \langle d^\dagger_{m\downarrow} d_{m'\uparrow}\rangle d^\dagger_{m\uparrow} d_{m'\downarrow}
    \right]
\end{align}
Where we shift the energy by
\begin{align}
    \Delta E &= -U \sum_m \langle n_{m\uparrow}\rangle \langle n_{m\downarrow}\rangle
    - U' \sum_{m\neq m'} \langle n_{m\uparrow}\rangle \langle n_{m'\downarrow}\rangle \\\nonumber
    &- (U'-J) \sum_{m < m',\sigma} \langle n_{m\sigma} \rangle \langle n_{m'\sigma} \rangle \\\nonumber
    &+ J\sum_{m\neq m'} \left[
        \langle d^\dagger_{m\uparrow} d_{m'\uparrow} \rangle 
        \langle d^\dagger_{m'\downarrow} d_{m\downarrow}\rangle
        -\langle d^\dagger_{m\uparrow} d_{m\downarrow}\rangle 
        \langle d^\dagger_{m'\downarrow} d_{m'\uparrow}\rangle
    \right] \\\nonumber
    &+ J\sum_{m\neq m'}\left[
        \langle d^\dagger_{m\uparrow} d_{m'\uparrow}\rangle 
        \langle d^\dagger_{m\downarrow} d_{m'\downarrow} \rangle
        - \langle d^\dagger_{m\uparrow} d_{m' \downarrow} \rangle
        \langle d^\dagger_{m\downarrow}d_{m'\uparrow}\rangle
    \right]
\end{align}
\section{Density matrix}
The density matrix for the grand canonical ensemble is
\begin{gather}
    \hat \rho = \frac{e^{-\beta(\hat H-\mu \hat N)}}{Z}
\end{gather}
with $Z = \text{Tr} ~e^{-\beta (H - \mu N)}$. The Hamiltonian is non-interacting so it can be written on the energy basis
\begin{gather}
    \hat H = \sum_n \varepsilon_n d_n^\dagger d_n 
\end{gather}
where the creation operator in this basis is defined by $d_n^{\dagger} = \sum_k \langle k|n\rangle c^\dagger_k$. Our fock space is defined by the occupation number in each energy state, for our fermion system this is either zero or one. We label the Fock state like $|\{n_k\}\rangle$. These are eigenstates of the Hamiltonian
\begin{align}
    \hat H|\{n_k\}\rangle = \sum_j \varepsilon_j \hat d_j^\dagger \hat d_j \left(\prod_k 
    \left(d_k^\dagger\right)^{n_k}\right)
\end{align}
For $j\neq k$ the commutation is zero and we can flip the two to normal order which results in zero. The only non zero is for when $j=k$. Therefore for a single fock state we can write
\begin{gather}
    \hat H|\{n_k\}\rangle = \sum_j \varepsilon_j n_j |\{n_k\}\rangle
\end{gather}
the same reasoning allows us to write
\begin{gather}
    \hat N |\{n_k\}\rangle = \sum_jn_j|\{n_k\}\rangle
\end{gather}
using these results we find the partition function
\begin{align}
    Z &= \sum_{n_0}\sum_{n_1}\ldots\sum_{n_N} \langle \{n_k\}|e^{-\beta(\hat H - \mu \hat N)}|\{n_k\}\rangle \\\nonumber
    &=\sum_{n_0}\sum_{n_1}\ldots\sum_{n_N} \langle \{n_k\}|e^{-\sum_j\beta(\varepsilon_j - \mu)n_j}|\{n_k\}\rangle \\\nonumber
    &= \sum_{n_0}\sum_{n_1}\ldots\sum_{n_N} \prod_j  \langle \{n_k\}|e^{-\beta(\varepsilon_j-\mu)n_j}|\{n_k\}\rangle
\end{align}
The sums out front go from 0, 1 since these are fermions. This means that the sums decouple and we can bring the product outside. Then we evaluate one sum from 0 to 1 and we get
\begin{gather}
    Z = \prod_j \left(1 + e^{-\beta(\varepsilon_j - \mu)}\right)
\end{gather}
Next we are interested in the density
\begin{align}
    \langle d^\dagger_m d_n \rangle &= \frac{\text{Tr}\left(e^{-\beta(\hat H - \mu \hat N)}\hat d_m^\dagger d_n\right)}{Z} \\\nonumber
    &= \frac{1}{Z} \sum_{n_0}\sum_{n_1}\ldots\sum_{n_N} \langle \{n_k\}|e^{-\beta(\hat H - \mu \hat N)}\hat d_m^\dagger d_n|\{n_k\}\rangle
\end{align}
Since our $\hat H$ and $\hat N$ are diagonal the density matrix is diagonal in this basis and $m = n$.
\begin{align}
    \langle d^\dagger_n d_n \rangle &=\frac{1}{Z}\sum_{n_0}\sum_{n_1}\ldots\sum_{n_N} \langle \{n_k\}|e^{-\beta(\hat H - \mu \hat N)}\hat d_n^\dagger d_n|\{n_k\}\rangle \\\nonumber
    &=\frac{1}{Z} \sum_{n_0}\sum_{n_1}\ldots\sum_{n_N} \langle \{n_k\}|e^{-\beta(\hat H - \mu \hat N)}n_n|\{n_k\}\rangle \\\nonumber
    &=\frac{1}{Z} \sum_{n_0}\sum_{n_1}\ldots\sum_{n_N} \langle \{n_k\}|e^{-\sum_j\beta(\varepsilon_j - \mu)n_j}n_n|\{n_k\}\rangle \\\nonumber
    &=\frac{1}{Z} \sum_{n_0}\sum_{n_1}\ldots\sum_{n_N} \langle \{n_k\}|n_n \prod_je^{-\beta(\varepsilon_j - \mu)n_j}|\{n_k\}\rangle
\end{align}
We can cancel out all $j\ne n$ terms from the partition function in the denominator. Hence we will get
\begin{align}
      \langle d^\dagger_n d_n \rangle &=\frac{\sum_{n=0}^1n_ne^{-\beta(\varepsilon_n - \mu)n_n}}{1+e^{-\beta(\varepsilon_n-\mu)}} \\\nonumber
      &= \frac{e^{-\beta(\varepsilon_n - \mu)}}{1+e^{-\beta(\varepsilon_n-\mu)}} \\\nonumber
      &= f(\varepsilon_n-\mu )
\end{align}
where $f$ is the fermi dirac distribution. We wish to write our density matrix in terms of the eigenstates. (diagonal operators d). To derive this
we can use bra-ket.
\begin{align}
    | c_a \rangle &= \sum_k |d_k\rangle \langle d_k | a \rangle \\\nonumber
    c^\dagger_a &= \sum_k \langle k | a \rangle d^\dagger_k
\end{align}
likewise
\begin{align}
    \langle c_b | &= \sum_q \langle c_q | d_q \rangle \langle d_q | \\\nonumber
    c_q &= \sum_q \langle b | q \rangle d_q
\end{align}
From the definition of the density matrix we can write
\begin{align}
    \rho_{ab} = \langle c_a^\dagger c_b\rangle &= \langle \left( \sum_k \langle k|a\rangle d_k^\dagger \right ) \left(\sum_q \langle b|q\rangle d_q \right) \rangle \\\nonumber
    &= \sum_{kq}\langle k|a\rangle\langle b | q\rangle \langle d_k^\dagger d_q\rangle \\\nonumber
    &= \sum_{kq} \langle b|q\rangle\langle k|a\rangle f(\varepsilon_k - \mu)\delta_{kq} \\\nonumber
    &= \sum_n \langle b|n\rangle\langle n|a\rangle f(\varepsilon_n -\mu)
\end{align}
Where $|n\rangle$ labels the eigenstates in our diagonal basis and $a,b$ labels the t2g,spin band. To implement
this in code it is nice to write this as an outer product recall
\begin{gather}
    \bm u_i \otimes \bm v_j =\bm u \cdot \bm v^T= \bm u_i \bm v_j
\end{gather}
where T is just the transpose not adjoint. For one energy slice $n$ elementwise we can write
\begin{align}
    \rho_{ab} = f_n \langle b | n \rangle \langle n | a \rangle = f_n v_b v_a^* = f_n v_a^* v_b
\end{align}
to get the index correct we write the outer product like
\begin{gather}
    \bm v_a \otimes \bm v_b
\end{gather}
which gives us the correct index structure but we need the conjugate
\begin{gather}
    \bm v_a^* \otimes \bm v_b
\end{gather}
hence in code we implement like
\begin{gather}
    \rho = \sum_n f_n \bm v_n^* \cdot \bm v_n^T
\end{gather}
which a standard matrix library can handle (I'm using eigen c++)
\end{document}